% claim.tex -- AI4Math claim of record (AI-generated, 2026-09-28; instantiated
% from harness/templates/claim/claim.tex per SOP 08b).
% Machine-readable theorem anchors are the "% LEAN:" lines; scripts/paper-lint.sh
% and the claim audit check them. All Lean listings are verbatim, byte-identical
% contiguous excerpts of frozen artifacts of tasks/20260928-comb08-sidemid:
%   statements  : formalized/01-comb08-sidemid-statements.lean (109 lines, sha256 5fe889d5...)
%   proof bodies: final/01-comb08-sidemid-proved.lean    (238 lines, sha256 19b5c390...)
% Line ranges used: statement 28-39, 41-74, 77-86, 89-108; final 80-115, 117-142,
% 144-159, 163-182 -- the same eight blocks the paper package verified byte-exact;
% bodies were injected by script (audit/inspect-tex.py --splice), never hand-copied.
% Verification: audit/listing-verify.txt (audit/inspect-tex.py --verify).
\documentclass[11pt]{article}

\usepackage[T1]{fontenc}
\usepackage[utf8]{inputenc}
\usepackage{amsmath,amssymb,amsthm}
\usepackage{graphicx}
\usepackage{hyperref}
\usepackage{xurl}
\setlength{\emergencystretch}{3em}
% lstlean.tex — Lean style for the listings package.
% Lineage: Jeremy Avigad (2015), adapted from Assia Mahboubi's SSR style;
% inspired by Olivier Verdier's unixode.sty for the Unicode literate table.
% No CTAN package or upstream repo exists; papers carry this file in their
% source package (cap set arXiv:1907.01449, sphere eversion arXiv:2210.07746).
% AI4Math adaptation (AI-generated, 2026-09-26): sorry/admit/sorryAx elevated
% to a dedicated error tier, matching the pipeline's 0-sorry constitution.
\usepackage{listings}
\usepackage{xcolor}

\definecolor{keywordcolor}{rgb}{0.7, 0.1, 0.1}   % red keywords
\definecolor{tacticcolor}{rgb}{0.0, 0.1, 0.6}    % blue tactics
\definecolor{commentcolor}{rgb}{0.4, 0.4, 0.4}   % grey comments
\definecolor{symbolcolor}{rgb}{0.0, 0.1, 0.6}    % blue symbols
\definecolor{sortcolor}{rgb}{0.1, 0.5, 0.1}      % green sorts
\definecolor{errorcolor}{rgb}{0.6, 0, 0}         % dark red: sorry/admit
\definecolor{stringcolor}{rgb}{0.5, 0.1, 0.5}    % purple strings

\lstdefinelanguage{lean}{
  % Anything between $ becomes LaTeX math mode
  mathescape=false,
  % Comments may or not include Latex commands
  texcl=false,
  % keywords, list taken from lean-syntax.txt
  morekeywords=[1]{import, prelude, protected, private, noncomputable,
    definition, meta, renaming, hiding, parameter, parameters, begin,
    constant, constants, lemma, variable, variables, theory, print,
    theorem, example, open, section, namespace, universe, universes, end,
    modifier, modifiers, using, export, exposing, axiom, inductive,
    coinductive, with, structure, class, instance, attribute, local, where,
    by, have, show, let, in, if, then, else, match, do, for, fun, assume,
    from, take, calc, include, omit, infix, infixl, infixr, notation,
    macro, syntax, elab, set_option, initialize, deriving, opaque, abbrev,
    mutual, partial, scoped, termination_by, decreasing_by},
  % Sorts
  morekeywords=[2]{Sort, Type, Prop},
  % tactics, list taken from lean-syntax.txt (tactic keywords)
  morekeywords=[3]{intro, intros, apply, exact, refine, rfl, simp, simp_all,
    simpa, simp_rw, rw, rewrite, erewrite, ring, ring_nf, omega, decide,
    norm_num, norm_cast, induction, cases, rcases, obtain, constructor,
    ext, funext, conv, congr, field_simp, linarith, nlinarith, positivity,
    tauto, trivial, aesop, use, by_cases, by_contra, push_neg, contrapose,
    symm, trans, assumption, contradiction, exfalso, specialize,
    generalize, revert, clear, rename_i, subst, unfold, delta, change,
    convert, split_ifs, interval_cases, fin_cases, zify, gcongr,
    nth_rewrite, suffices, generalizing, at, only, native_decide},
  % warnings - constitution tier: must never survive into a final proof
  morekeywords=[4]{sorry, admit, sorryAx},
  literate=
    {α}{{\ensuremath{\alpha}}}1
    {β}{{\ensuremath{\beta}}}1
    {γ}{{\ensuremath{\gamma}}}1
    {δ}{{\ensuremath{\delta}}}1
    {ε}{{\ensuremath{\varepsilon}}}1
    {ζ}{{\ensuremath{\zeta}}}1
    {η}{{\ensuremath{\eta}}}1
    {θ}{{\ensuremath{\theta}}}1
    {ι}{{\ensuremath{\iota}}}1
    {κ}{{\ensuremath{\kappa}}}1
    {λ}{{\ensuremath{\lambda}}}1
    {μ}{{\ensuremath{\mu}}}1
    {ν}{{\ensuremath{\nu}}}1
    {ξ}{{\ensuremath{\xi}}}1
    {π}{{\ensuremath{\pi}}}1
    {ρ}{{\ensuremath{\rho}}}1
    {σ}{{\ensuremath{\sigma}}}1
    {τ}{{\ensuremath{\tau}}}1
    {υ}{{\ensuremath{\upsilon}}}1
    {φ}{{\ensuremath{\varphi}}}1
    {χ}{{\ensuremath{\chi}}}1
    {ψ}{{\ensuremath{\psi}}}1
    {ω}{{\ensuremath{\omega}}}1
    {Γ}{{\ensuremath{\Gamma}}}1
    {Δ}{{\ensuremath{\Delta}}}1
    {Θ}{{\ensuremath{\Theta}}}1
    {Λ}{{\ensuremath{\Lambda}}}1
    {Ξ}{{\ensuremath{\Xi}}}1
    {Π}{{\ensuremath{\Pi}}}1
    {Σ}{{\ensuremath{\Sigma}}}1
    {Φ}{{\ensuremath{\Phi}}}1
    {Ψ}{{\ensuremath{\Psi}}}1
    {Ω}{{\ensuremath{\Omega}}}1
    {ℵ}{{\ensuremath{\aleph}}}1
    {∀}{{\ensuremath{\forall}}}1
    {∃!}{{\ensuremath{\exists}!}}1
    {∃}{{\ensuremath{\exists}}}1
    {∈}{{\ensuremath{\in}}}1
    {∉}{{\ensuremath{\notin}}}1
    {∘}{{\ensuremath{\circ}}}1
    {∩}{{\ensuremath{\cap}}}1
    {∪}{{\ensuremath{\cup}}}1
    {⊆}{{\ensuremath{\subseteq}}}1
    {⊂}{{\ensuremath{\subset}}}1
    {⊇}{{\ensuremath{\supseteq}}}1
    {⊃}{{\ensuremath{\supset}}}1
    {⊄}{{\ensuremath{\nsubseteq}}}1
    {∅}{{\ensuremath{\emptyset}}}1
    {∧}{{\ensuremath{\wedge}}}1
    {∨}{{\ensuremath{\vee}}}1
    {¬}{{\ensuremath{\neg}}}1
    {⊤}{{\ensuremath{\top}}}1
    {⊥}{{\ensuremath{\bot}}}1
    {↔}{{\ensuremath{\leftrightarrow}}}1
    {→}{{\ensuremath{\rightarrow}}}1
    {←}{{\ensuremath{\leftarrow}}}1
    {↑}{{\ensuremath{\uparrow}}}1
    {⇒}{{\ensuremath{\Rightarrow}}}1
    {⇐}{{\ensuremath{\Leftarrow}}}1
    {⟹}{{\ensuremath{\Longrightarrow}}}1
    {↦}{{\ensuremath{\mapsto}}}1
    {≤}{{\ensuremath{\leq}}}1
    {≥}{{\ensuremath{\geq}}}1
    {≠}{{\ensuremath{\neq}}}1
    {≈}{{\ensuremath{\approx}}}1
    {≡}{{\ensuremath{\equiv}}}1
    {≃}{{\ensuremath{\simeq}}}1
    {≅}{{\ensuremath{\cong}}}1
    {×}{{\ensuremath{\times}}}1
    {·}{{\ensuremath{\cdot}}}1
    {±}{{\ensuremath{\pm}}}1
    {⊕}{{\ensuremath{\oplus}}}1
    {⊗}{{\ensuremath{\otimes}}}1
    {⋆}{{\ensuremath{\star}}}1
    {▸}{{\ensuremath{\blacktriangleright}}}1
    {⟨}{{\ensuremath{\langle}}}1
    {⟩}{{\ensuremath{\rangle}}}1
    {⦃}{{\ensuremath{\{\!\mid}}}1
    {⦄}{{\ensuremath{\mid\!\}}}}1
    {⟦}{{\ensuremath{[\![}}}1
    {⟧}{{\ensuremath{]\!]}}}1
    {⌊}{{\ensuremath{\lfloor}}}1
    {⌋}{{\ensuremath{\rfloor}}}1
    {⌈}{{\ensuremath{\lceil}}}1
    {⌉}{{\ensuremath{\rceil}}}1
    {√}{{\ensuremath{\surd}}}1
    {∣}{{\ensuremath{\mid}}}1
    {∤}{{\ensuremath{\nmid}}}1
    {∎}{{\ensuremath{\blacksquare}}}1
    {⋯}{{\ensuremath{\cdots}}}1
    {…}{{\ensuremath{\ldots}}}1
    {∑}{{\ensuremath{\sum}}}1
    {∏}{{\ensuremath{\prod}}}1
    {⁻¹}{{\ensuremath{^{-1}}}}1
    {¹}{{\ensuremath{^1}}}1
    {²}{{\ensuremath{^2}}}1
    {³}{{\ensuremath{^3}}}1
    {ⁿ}{{\ensuremath{^n}}}1
    {ᵐ}{{\ensuremath{^m}}}1
    {₀}{{\ensuremath{_0}}}1
    {₁}{{\ensuremath{_1}}}1
    {₂}{{\ensuremath{_2}}}1
    {₃}{{\ensuremath{_3}}}1
    {₄}{{\ensuremath{_4}}}1
    {₅}{{\ensuremath{_5}}}1
    {ₙ}{{\ensuremath{_n}}}1
    {ₘ}{{\ensuremath{_m}}}1
    {ₖ}{{\ensuremath{_k}}}1
    {ᵢ}{{\ensuremath{_i}}}1
    {ⱼ}{{\ensuremath{_j}}}1
    {ℤ}{{\ensuremath{\mathbb{Z}}}}1
    {ℕ}{{\ensuremath{\mathbb{N}}}}1
    {ℚ}{{\ensuremath{\mathbb{Q}}}}1
    {ℝ}{{\ensuremath{\mathbb{R}}}}1
    {ℂ}{{\ensuremath{\mathbb{C}}}}1
    {𝔽}{{\ensuremath{\mathbb{F}}}}1
    {𝒫}{{\ensuremath{\mathcal{P}}}}1
    {∗}{{\ensuremath{\ast}}}1
    {•}{{\ensuremath{\bullet}}}1,
  % Comments
  morecomment=[l]{--},
  morecomment=[s]{/-}{-/},
  morestring=[b]",
  stringstyle=\color{stringcolor},
  keepspaces=true,
  keywordstyle=[1]\color{keywordcolor}\bfseries,
  keywordstyle=[2]\color{sortcolor}\bfseries,
  keywordstyle=[3]\color{tacticcolor}\bfseries,
  keywordstyle=[4]\color{errorcolor}\bfseries\underline,
  escapeinside={(*@}{@*)},
  columns=fullflexible,
  basicstyle=\ttfamily\small,
  commentstyle=\color{commentcolor}\itshape,
  breaklines=true,
  showstringspaces=false,
  frame=single,
  framesep=2mm,
  xleftmargin=4mm,
  numbers=left,
  numberstyle=\tiny\color{commentcolor},
  numbersep=6pt,
}

% Inline Lean code in prose: \lean{Int.ModEq.pow}
\newcommand{\lean}[1]{\lstinline[language=lean,basicstyle=\ttfamily\normalsize]|#1|}

\lstset{language=lean}

% --- local rendering patch (2026-09-28) ---
% tectonic/XeTeX does not apply the listings literate table to multi-byte
% characters, so the Unicode set used by the listings of THIS note is mapped to
% math glyphs as active characters via the standard \lccode/\lowercase idiom
% (ASCII-only source, no extra packages). Same block as the paper package of
% this result: the template block plus the four code points that occur in the
% frozen comb08-sidemid sources but are absent from the template patch
% (2212 21A6 2208 27FA). Restricted to symbols outside xeCJK's own punctuation
% table: 00B7, 2014, 3002 and the fullwidth forms also occur in the frozen
% sources but are special punctuation that xeCJK sets up at load time -- making
% them active first aborts the run with "Control sequence ... already defined"
% (observed 2026-09-27/28); those route through xeCJK + SimSun instead.
% Rendering only: listing source bytes are untouched, so the byte-exactness
% checks are unaffected, and the \ifdefined guard leaves the pdflatex path
% identical to the template.
\ifdefined\XeTeXversion
  {\lccode`\~="2022 \lowercase{\global\catcode`~=\active \global\def~{\ensuremath{\bullet}}}}
  {\lccode`\~="2115 \lowercase{\global\catcode`~=\active \global\def~{\ensuremath{\mathbb{N}}}}}
  {\lccode`\~="2124 \lowercase{\global\catcode`~=\active \global\def~{\ensuremath{\mathbb{Z}}}}}
  {\lccode`\~="2190 \lowercase{\global\catcode`~=\active \global\def~{\ensuremath{\leftarrow}}}}
  {\lccode`\~="2192 \lowercase{\global\catcode`~=\active \global\def~{\ensuremath{\rightarrow}}}}
  {\lccode`\~="2194 \lowercase{\global\catcode`~=\active \global\def~{\ensuremath{\leftrightarrow}}}}
  {\lccode`\~="2200 \lowercase{\global\catcode`~=\active \global\def~{\ensuremath{\forall}}}}
  {\lccode`\~="2203 \lowercase{\global\catcode`~=\active \global\def~{\ensuremath{\exists}}}}
  {\lccode`\~="2227 \lowercase{\global\catcode`~=\active \global\def~{\ensuremath{\wedge}}}}
  {\lccode`\~="2228 \lowercase{\global\catcode`~=\active \global\def~{\ensuremath{\vee}}}}
  {\lccode`\~="2261 \lowercase{\global\catcode`~=\active \global\def~{\ensuremath{\equiv}}}}
  {\lccode`\~="2264 \lowercase{\global\catcode`~=\active \global\def~{\ensuremath{\leq}}}}
  {\lccode`\~="2265 \lowercase{\global\catcode`~=\active \global\def~{\ensuremath{\geq}}}}
  {\lccode`\~="22A2 \lowercase{\global\catcode`~=\active \global\def~{\ensuremath{\vdash}}}}
  {\lccode`\~="27E8 \lowercase{\global\catcode`~=\active \global\def~{\ensuremath{\langle}}}}
  {\lccode`\~="27E9 \lowercase{\global\catcode`~=\active \global\def~{\ensuremath{\rangle}}}}
  {\lccode`\~="2212 \lowercase{\global\catcode`~=\active \global\def~{\ensuremath{-}}}}
  {\lccode`\~="21A6 \lowercase{\global\catcode`~=\active \global\def~{\ensuremath{\mapsto}}}}
  {\lccode`\~="2208 \lowercase{\global\catcode`~=\active \global\def~{\ensuremath{\in}}}}
  {\lccode`\~="27FA \lowercase{\global\catcode`~=\active \global\def~{\ensuremath{\Longleftrightarrow}}}}
\fi
\ifdefined\XeTeXversion
  \usepackage{xeCJK}
  \setCJKmainfont{SimSun}
\fi

\newtheorem{theorem}{Theorem}
\theoremstyle{definition}
\newtheorem{conjecture}{Conjecture}

\title{Tiling-count sequence of the $3\times n$ rectangle with a side-middle
cell removed: interleave identity and parity/mod-4 laws\\ (Lean 4
machine-checked claim of record)}
\author{Aurora0134\thanks{Contact: via the GitHub account Aurora0134; ORCID
placeholder, to be filled in by the author at Zenodo upload. No personal
information is carried in this note.}}
\date{September 28, 2026}

\begin{document}
\maketitle

\begin{abstract}
Let $a(n)$ count the domino tilings of a $3\times n$ rectangle with the
middle cell of its last column removed. The sequence
$1,8,46,295,1814,11306,\dots$ is defined by six seeds and the order-six sparse
recurrence $a(n+6)=4a(n+5)+14a(n+4)-10a(n+2)+a(n)$. We state and
machine-check, in Lean 4 (v4.34.0, mathlib v4.34.0), three theorems: the main
sequence is the interleaving of its odd- and even-indexed subsequences, which
satisfy one and the same order-six recurrence; $a(n)$ is odd if and only if
$n\equiv 1\pmod 3$; and $a(n)\bmod 4$ depends only on $n\bmod 12$ with
pattern $(1,0,2,3,2,2,1,2,0,3,0,0)$. Every proof compiles with no
\texttt{sorry} and \texttt{\#print axioms} reports only \{propext,
Quot.sound\}. Novelty adjudication: no occupying record found in three review
rounds; value tier \emph{new sequence} --- the recurrence's order and
coefficients are shared with A033506 and no novelty is claimed for them. The
combinatorial identification is not kernel-proved and is stated as a
conjecture. This note is a priority claim of record, not a full paper.
\end{abstract}

\section{Introduction}

The object of this note is the tiling-count sequence of the $3\times n$
rectangle with a side-middle cell removed: tiled by $1\times 2$ dominoes, the
number of tilings is $a(n)$, and the sequence begins
$1,8,46,295,1814,11306,70161,\dots$ for $n\ge 1$ (the matching-polynomial
setting \cite{godsil-gutman}). Three pure-sequence theorems about it --- an
interleave identity and two modular laws --- were produced by the AI4Math
pipeline, in which LLM workers generate statements and proof scripts while the
Lean 4 kernel is the sole adjudicator at every gate: statements are frozen by
hash before proving, and a separate audit re-computes the symmetric
difference, scans the whole file for \texttt{sorry}/\texttt{admit} including
comments, and re-prints the axiom dependencies \cite{lean4,mathlib}.
Verification strength is stated below verbatim. This note is a priority claim
of record deposited so that a dated, citable public record exists ahead of any
narrative paper; the fuller paper for the same theorems is prepared separately
and will reference this deposit. No literature review is attempted here.

\section{Statement}\label{sec:statement}

Indexing conventions, locked on the source task: $0$-based in the Lean layer,
$1$-based in the counting reading, through $A(n)=a(n+1)$, $O(k)=a(2k+1)$ and
$E(k)=a(2k+2)$; the Lean sequences \lean{asid}, \lean{osid} and \lean{esid}
are $A$, $O$ and $E$. The statements below are about sequences defined by
literal initial values and integer-coefficient recurrences only; no graph
object appears in the frozen statements. The main sequence is defined by six
initial values and the order-six sparse recurrence
\[
  a(n+6)=4a(n+5)+14a(n+4)-10a(n+2)+a(n),
\]
whose signature $(4,14,0,-10,0,1)$ is the one registered for the undefected
$3\times n$ grid, A033506 \cite{oeis-a033506}; nothing below depends on, and
nothing below claims anything about, that sharing.

\begin{lstlisting}[caption={The main sequence \texttt{asid}, verbatim from the frozen statement file \texttt{formalized}\allowbreak/\texttt{01-comb08-sidemid-statements.lean} (lines 28--39, sha256 \texttt{5fe889d5}).},label={lst:asid}]
/-- **序列 asid（主列，六阶稀疏递推）**。
初值 a(0..5) = 1, 8, 46, 295, 1814, 11306；
递推 a(n+6) = 4·a(n+5) + 14·a(n+4) − 10·a(n+2) + a(n)（n ≥ 0；
递推签名 (4,14,0,−10,0,1)，系数为 0 的滞后项按主列 def 口径省略不写）。 -/
def asid : ℕ → ℤ
  | 0 => 1
  | 1 => 8
  | 2 => 46
  | 3 => 295
  | 4 => 1814
  | 5 => 11306
  | n + 6 => 4 * asid (n + 5) + 14 * asid (n + 4) - 10 * asid (n + 2) + asid n
\end{lstlisting}

% LEAN: tasks/20260928-comb08-sidemid :: asid_interleave grade=完全证明
\begin{theorem}[interleaving]\label{thm:interleave}
For every $k\in\mathbb{N}$,
\[
  A(2k)=O(k)\qquad\text{and}\qquad A(2k+1)=E(k).
\]
In the $1$-based reading this is $a(2k+1)=o(k)$ and $a(2k+2)=e(k)$: the main
sequence is the interleaving of its odd-indexed and even-indexed
subsequences, and those two subsequences satisfy one and the same order-six
recurrence.
\end{theorem}

\begin{lstlisting}[caption={Sequences \texttt{osid}, \texttt{esid} and the statement of Theorem~\ref{thm:interleave}, verbatim from the frozen statement file (lines 41--74, sha256 \texttt{5fe889d5}). The \texttt{:= by} at the end of the theorem line opens the proof body, which is given in Listing~\ref{lst:proof-t1}; the \texttt{sorry} placeholder of the frozen file is not part of the excerpt.},label={lst:stmt-t1}]
/-- **序列 osid（1-based 奇位子列，六阶递推）**。
初值 o(0..5) = 1, 46, 1814, 70161, 2708104, 104507480
（= asid 在下标 0,2,4,6,8,10 处之值，即 a(1),a(3),a(5),a(7),a(9),a(11)）；
递推 o(k+6) = 44·o(k+5) − 216·o(k+4) + 282·o(k+3) − 128·o(k+2) + 20·o(k+1) − o(k)（k ≥ 0）。 -/
def osid : ℕ → ℤ
  | 0 => 1
  | 1 => 46
  | 2 => 1814
  | 3 => 70161
  | 4 => 2708104
  | 5 => 104507480
  | k + 6 =>
      44 * osid (k + 5) - 216 * osid (k + 4) + 282 * osid (k + 3)
        - 128 * osid (k + 2) + 20 * osid (k + 1) - osid k

/-- **序列 esid（1-based 偶位子列，六阶递推）**。
初值 e(0..5) = 8, 295, 11306, 435986, 16823455, 649209736
（= asid 在下标 1,3,5,7,9,11 处之值，即 a(2),a(4),a(6),a(8),a(10),a(12)）；
递推 e(k+6) = 44·e(k+5) − 216·e(k+4) + 282·e(k+3) − 128·e(k+2) + 20·e(k+1) − e(k)（k ≥ 0）。 -/
def esid : ℕ → ℤ
  | 0 => 8
  | 1 => 295
  | 2 => 11306
  | 3 => 435986
  | 4 => 16823455
  | 5 => 649209736
  | k + 6 =>
      44 * esid (k + 5) - 216 * esid (k + 4) + 282 * esid (k + 3)
        - 128 * esid (k + 2) + 20 * esid (k + 1) - esid k

/-- **T1（主攻）**：主列六阶递推 = 两条同系数六阶子列的交织——
对一切 k，asid (2k) = osid k 且 asid (2k+1) = esid k。 -/
theorem asid_interleave (k : ℕ) :
    asid (2 * k) = osid k ∧ asid (2 * k + 1) = esid k := by
\end{lstlisting}

% LEAN: tasks/20260928-comb08-sidemid :: asid_mod2_period3 grade=完全证明
\begin{theorem}[modulo $2$ pattern]\label{thm:mod2}
For every $n\in\mathbb{N}$ and every $r$ with $n\bmod 3=r$,
\[
  A(n)\bmod 2=P_3(r),\qquad P_3(0)=1,\quad P_3(r)=0\ \text{for }r\neq 0.
\]
In the $1$-based convention: $a(n)$ is odd if and only if $n\equiv
1\pmod 3$.
\end{theorem}

\begin{lstlisting}[caption={Pattern \texttt{pat3} and the statement of Theorem~\ref{thm:mod2}, verbatim from the frozen statement file (lines 77--86, sha256 \texttt{5fe889d5}).},label={lst:stmt-t2}]
/-- **余数样式 pat3**：0↦1，其余↦0
（mod 2 余数 0-based 口径：asid n 为奇 ⟺ n ≡ 0 (mod 3)）。 -/
def pat3 : ℕ → ℤ
  | 0 => 1
  | _ => 0

/-- **T2（伴随）**：asid 的 mod 2 余数仅依赖 n mod 3，样式为 pat3
（即 asid n 为奇 ⟺ n ≡ 0 (mod 3)；1-based 口径 = a(n) 为奇 ⟺ n ≡ 1 (mod 3)）。 -/
theorem asid_mod2_period3 (n r : ℕ) (hr : n % 3 = r) :
    asid n % 2 = pat3 r := by
\end{lstlisting}

% LEAN: tasks/20260928-comb08-sidemid :: asid_mod4_period12 grade=完全证明
\begin{theorem}[modulo $4$ pattern]\label{thm:mod4}
For every $n\in\mathbb{N}$ and every $r$ with $n\bmod 12=r$,
\[
  A(n)\bmod 4=P_{12}(r),
\]
where $P_{12}$ on the residues $0,\dots,11$ is
$(1,0,2,3,2,2,1,2,0,3,0,0)$. In the $1$-based convention: $a(n)\bmod 4$
depends only on $n\bmod 12$, with the same pattern along the residues
$1,\dots,12$. The theorem implies Theorem~\ref{thm:mod2}: the odd values
occur exactly at the residues $0,3,6,9$ modulo $12$, that is, at
$n\equiv 0\pmod 3$ in the $0$-based reading.
\end{theorem}

\begin{lstlisting}[caption={Pattern \texttt{pat12} and the statement of Theorem~\ref{thm:mod4}, verbatim from the frozen statement file (lines 89--108, sha256 \texttt{5fe889d5}).},label={lst:stmt-t3}]
/-- **余数样式 pat12**（12 留数类 ↦ mod 4 余数）：
0↦1, 1↦0, 2↦2, 3↦3, 4↦2, 5↦2, 6↦1, 7↦2, 8↦0, 9↦3, 10↦0, 11↦0。 -/
def pat12 : ℕ → ℤ
  | 0 => 1
  | 1 => 0
  | 2 => 2
  | 3 => 3
  | 4 => 2
  | 5 => 2
  | 6 => 1
  | 7 => 2
  | 8 => 0
  | 9 => 3
  | 10 => 0
  | _ => 0

/-- **T3（伴随）**：asid 的 mod 4 余数仅依赖 n mod 12，样式为 pat12
（蕴含 T2：奇值恰在 n%12 ∈ {0,3,6,9} ⟺ n ≡ 0 (mod 3)）。 -/
theorem asid_mod4_period12 (n r : ℕ) (hr : n % 12 = r) :
    asid n % 4 = pat12 r := by
\end{lstlisting}

\begin{conjecture}[combinatorial bridge, not proved here]\label{conj:bridge}
For every $n\ge 1$, $A(n-1)=a(n)$, where $a$ counts the matchings (including
the empty one) of the graph obtained from the $3\times n$ grid by deleting
the middle vertex of the last column.
\end{conjecture}

\section{Proof}\label{sec:proof}

All proofs share one skeleton: strong induction on the index
(\lean{Nat.strong_induction_on}), the seed window discharged as literal cases
by \texttt{interval\_cases} plus \texttt{decide}, one definitional unfolding
of the recurrence in the inductive step (\texttt{simp only [asid]} followed
by \texttt{omega}), and residue-class case analysis by \texttt{mod\_cases}
where the goal branches on a residue. The auxiliary twelve-step even-lag
identity \lean{asid_evenlag12} (Listing~\ref{lst:proof-aux}) is what makes
the term-by-term matching in the interleave proof possible; its complete
proof, together with the one-step unfolding \lean{asid_shift} it builds on, is
reproduced verbatim below. The complete proofs of
Theorems~\ref{thm:interleave} and \ref{thm:mod2} follow verbatim; for
Theorem~\ref{thm:mod4} the statement, the induction skeleton and the first of
the twelve \texttt{mod\_cases} branches are excerpted --- the remaining eleven
branches are identical in shape (lines 183--237 of the final file) and the
complete $238$-line file ships in this deposit under \texttt{proofs/}.

\begin{lstlisting}[caption={The auxiliary lemmas \texttt{asid\_shift} and \texttt{asid\_evenlag12} with complete proofs, verbatim from \texttt{final}\allowbreak/\texttt{01-comb08-sidemid-proved.lean} (lines 80--115, sha256 \texttt{19b5c390}). These two declarations are proof-body scaffolding added after the freeze and therefore do not occur in the frozen statement file.},label={lst:proof-aux}]
/-- 辅助：asid 单步展开（t+6 形）。 -/
theorem asid_shift (t : ℕ) :
    asid (t + 6) = 4 * asid (t + 5) + 14 * asid (t + 4) - 10 * asid (t + 2) + asid t := by
  simp only [asid]

/-- 辅助：12 阶偶滞后恒等式 asid(m+12) = 44·asid(m+10) − 216·asid(m+8) + 282·asid(m+6)
− 128·asid(m+4) + 20·asid(m+2) − asid(m)。全级联展开到基 asid(m..m+5) 后 ring 闭合。 -/
theorem asid_evenlag12 (m : ℕ) :
    asid (m + 12) = 44 * asid (m + 10) - 216 * asid (m + 8) + 282 * asid (m + 6)
      - 128 * asid (m + 4) + 20 * asid (m + 2) - asid m := by
  rw [show m + 12 = (m + 6) + 6 by ring, asid_shift (m + 6),
      show m + 6 + 5 = m + 11 by ring,
      show m + 6 + 4 = m + 10 by ring,
      show m + 6 + 2 = m + 8 by ring]
  rw [show m + 11 = (m + 5) + 6 by ring, asid_shift (m + 5),
      show m + 5 + 5 = m + 10 by ring,
      show m + 5 + 4 = m + 9 by ring,
      show m + 5 + 2 = m + 7 by ring]
  rw [show m + 10 = (m + 4) + 6 by ring, asid_shift (m + 4),
      show m + 4 + 5 = m + 9 by ring,
      show m + 4 + 4 = m + 8 by ring,
      show m + 4 + 2 = m + 6 by ring]
  rw [show m + 9 = (m + 3) + 6 by ring, asid_shift (m + 3),
      show m + 3 + 5 = m + 8 by ring,
      show m + 3 + 4 = m + 7 by ring,
      show m + 3 + 2 = m + 5 by ring]
  rw [show m + 8 = (m + 2) + 6 by ring, asid_shift (m + 2),
      show m + 2 + 5 = m + 7 by ring,
      show m + 2 + 4 = m + 6 by ring,
      show m + 2 + 2 = m + 4 by ring]
  rw [show m + 7 = (m + 1) + 6 by ring, asid_shift (m + 1),
      show m + 1 + 5 = m + 6 by ring,
      show m + 1 + 4 = m + 5 by ring,
      show m + 1 + 2 = m + 3 by ring]
  rw [asid_shift m]
  ring
\end{lstlisting}

\begin{lstlisting}[caption={Complete proof of Theorem~\ref{thm:interleave}, verbatim from \texttt{final}\allowbreak/\texttt{01-comb08-sidemid-proved.lean} (lines 117--142, sha256 \texttt{19b5c390}).},label={lst:proof-t1}]
/-- **T1（主攻）**：主列六阶递推 = 两条同系数六阶子列的交织——
对一切 k，asid (2k) = osid k 且 asid (2k+1) = esid k。 -/
theorem asid_interleave (k : ℕ) :
    asid (2 * k) = osid k ∧ asid (2 * k + 1) = esid k := by
  induction k using Nat.strong_induction_on with
  | h k ih =>
    by_cases hk : k < 6
    · interval_cases k <;> decide
    · obtain ⟨j, rfl⟩ : ∃ j, k = j + 6 := ⟨k - 6, by omega⟩
      constructor
      · rw [show 2 * (j + 6) = (2 * j) + 12 by ring, asid_evenlag12 (2 * j)]
        simp only [osid]
        rw [show 2 * j + 10 = 2 * (j + 5) by ring, (ih (j + 5) (by omega)).1,
            show 2 * j + 8 = 2 * (j + 4) by ring, (ih (j + 4) (by omega)).1,
            show 2 * j + 6 = 2 * (j + 3) by ring, (ih (j + 3) (by omega)).1,
            show 2 * j + 4 = 2 * (j + 2) by ring, (ih (j + 2) (by omega)).1,
            show 2 * j + 2 = 2 * (j + 1) by ring, (ih (j + 1) (by omega)).1,
            (ih j (by omega)).1]
      · rw [show 2 * (j + 6) + 1 = (2 * j + 1) + 12 by ring, asid_evenlag12 (2 * j + 1)]
        simp only [esid]
        rw [show 2 * j + 1 + 10 = 2 * (j + 5) + 1 by ring, (ih (j + 5) (by omega)).2,
            show 2 * j + 1 + 8 = 2 * (j + 4) + 1 by ring, (ih (j + 4) (by omega)).2,
            show 2 * j + 1 + 6 = 2 * (j + 3) + 1 by ring, (ih (j + 3) (by omega)).2,
            show 2 * j + 1 + 4 = 2 * (j + 2) + 1 by ring, (ih (j + 2) (by omega)).2,
            show 2 * j + 1 + 2 = 2 * (j + 1) + 1 by ring, (ih (j + 1) (by omega)).2,
            (ih j (by omega)).2]
\end{lstlisting}

\begin{lstlisting}[caption={Complete proof of Theorem~\ref{thm:mod2}, verbatim from \texttt{final}\allowbreak/\texttt{01-comb08-sidemid-proved.lean} (lines 144--159, sha256 \texttt{19b5c390}).},label={lst:proof-t2}]
/-- **T2（伴随）**：asid 的 mod 2 余数仅依赖 n mod 3，样式为 pat3
（即 asid n 为奇 ⟺ n ≡ 0 (mod 3)；1-based 口径 = a(n) 为奇 ⟺ n ≡ 1 (mod 3)）。 -/
theorem asid_mod2_period3 (n r : ℕ) (hr : n % 3 = r) :
    asid n % 2 = pat3 r := by
  have main : ∀ n : ℕ, asid n % 2 = pat3 (n % 3) := by
    intro n
    induction n using Nat.strong_induction_on with
    | h n ih =>
      by_cases hn : n < 6
      · interval_cases n <;> decide
      · obtain ⟨j, rfl⟩ : ∃ j, n = j + 6 := ⟨n - 6, by omega⟩
        have h1 : asid (j + 6) % 2 = asid j % 2 := by
          simp only [asid]
          omega
        rw [h1, ih j (by omega), show (j + 6) % 3 = j % 3 by omega]
  rw [← hr]; exact main n
\end{lstlisting}

\begin{lstlisting}[caption={Theorem~\ref{thm:mod4}: statement, induction skeleton and the first of the twelve \texttt{mod\_cases} branches, verbatim from \texttt{final}\allowbreak/\texttt{01-comb08-sidemid-proved.lean} (lines 163--182, sha256 \texttt{19b5c390}). The remaining eleven branches are identical in shape (lines 183--237) and ship in this deposit under \texttt{proofs/}.},label={lst:proof-t3}]
theorem asid_mod4_period12 (n r : ℕ) (hr : n % 12 = r) :
    asid n % 4 = pat12 r := by
  have main : ∀ n : ℕ, asid n % 4 = pat12 (n % 12) := by
    intro n
    induction n using Nat.strong_induction_on with
    | h n ih =>
      by_cases hn : n < 6
      · interval_cases n <;> decide
      · obtain ⟨j, rfl⟩ : ∃ j, n = j + 6 := ⟨n - 6, by omega⟩
        have h1 : asid (j + 6) % 4
            = (2 * (asid (j + 4) % 4) + 2 * (asid (j + 2) % 4) + asid j % 4) % 4 := by
          simp only [asid]
          omega
        rw [h1, ih (j + 4) (by omega), ih (j + 2) (by omega), ih j (by omega)]
        mod_cases hj : j % 12
        · have hj0 : j % 12 = 0 := by simpa [Nat.ModEq] using hj
          have h2 : (j + 2) % 12 = 2 := by omega
          have h4 : (j + 4) % 12 = 4 := by omega
          have h6 : (j + 6) % 12 = 6 := by omega
          rw [hj0, h2, h4, h6]; decide
\end{lstlisting}

\section{Verification evidence}

\begin{itemize}
\item Toolchain: Lean \texttt{v4.34.0}
(\texttt{leanprover/lean4:v4.34.0}), mathlib4 \texttt{v4.34.0}, revision
\texttt{5ed29652} \cite{lean4,mathlib,mathlib-v434}. Reproduction: with that
pin, \texttt{lake env lean} on \texttt{proofs/01-comb08-sidemid-proved.lean}
compiles the file with exit code $0$; the strict re-compile of the merged
final file measured $29$\,s wall clock in the task ledger.
\item Statement integrity: the frozen statement file
\texttt{formalized}\allowbreak/\texttt{01-comb08-sidemid-statements.lean}
($109$ lines) has sha256
\texttt{5fe889d5}\allowbreak\texttt{ed92711a}\allowbreak\texttt{86f23d98}\allowbreak\texttt{27800dd7}\allowbreak\texttt{d2c677a4}\allowbreak\texttt{40664f0d}\allowbreak\texttt{81dbdb5e}\allowbreak\texttt{9eb39e67};
the final file
\texttt{final}\allowbreak/\texttt{01-comb08-sidemid-proved.lean}
($238$ lines) has sha256
\texttt{19b5c390}\allowbreak\texttt{8592a529}\allowbreak\texttt{491e0b5b}\allowbreak\texttt{40b6b261}\allowbreak\texttt{66019d92}\allowbreak\texttt{820447b4}\allowbreak\texttt{1a70ada5}\allowbreak\texttt{10a20821}.
The final file was compared with the frozen file by symmetric difference over
the declaration headers: all $8$ headers byte-identical, zero drift; the only
differences are the three proof bodies replacing the \texttt{sorry}
placeholders and the two auxiliary lemmas added inside the proof bodies.
\item Kernel check: compiles with $0$ \texttt{sorry} (whole-file scan
including comments, case-insensitive, also for \texttt{admit} and
\texttt{sorryAx}); \texttt{\#print axioms} output, verbatim (kept literal;
the smaller monospace width is typesetting only):
\begin{lstlisting}[basicstyle=\ttfamily\footnotesize,frame=none]
'asid_interleave' depends on axioms: [propext, Quot.sound]
'asid_mod2_period3' depends on axioms: [propext, Quot.sound]
'asid_mod4_period12' depends on axioms: [propext, Quot.sound]
\end{lstlisting}
The auxiliary lemmas, probe-verified in the same strict pass, depend on
\texttt{[propext, Quot.sound]} (\lean{asid_evenlag12}) and \texttt{[propext]}
(\lean{asid_shift}); \texttt{sorryAx} and \texttt{Classical.choice} never
appear.
\item Grading by the audit department: all three theorems \emph{complete
proof}; lowest grade reached for the case = \emph{complete proof}; no
scaffolding route was used.
\item Budgets, copied from the source task ledger's final machine
reconciliation: agent-run sampling $14$ of a permitted $20$;
\texttt{llm-call} $1$ of a permitted $6$; kernel compiles $26$ of a permitted
$32$; LaTeX rounds $9$ of a permitted $10$.
\item Audit chain (originals in this deposit under \texttt{audit/}): final
adjudication \texttt{final-comb08-sidemid.txt} (four independent checks), the
three axiom probe outputs
\texttt{axioms-asid\_}\allowbreak\texttt{interleave.txt},
\texttt{axioms-asid\_}\allowbreak\texttt{mod2\_period3.txt} and
\texttt{axioms-asid\_}\allowbreak\texttt{mod4\_period12.txt}, and this note's
listing verification \texttt{listing-verify.txt}. The semantic-fidelity check
(roundtrip re-translation) of the underlying task ran on the same model node
as its formalisation; that independence weakness is mitigated, not removed,
by kernel adjudication and manual review.
\end{itemize}

\section{Novelty evidence}\label{sec:novelty}

Adjudication copied from the pipeline's exclusivity review (C9 gate), which
ran in three rounds around this task and re-derived the data independently
rather than trusting the candidate card. Verdict: \textbf{no occupying record
found} (three rounds consistent); value tier, as corrected by the third
round: \textbf{new sequence} --- the sequence is unoccupied, while its
recurrence is \emph{not} new. The full query logs, zero-hit entries included,
are in this deposit at \texttt{audit/c9-record.md} (verbatim copies of the
review records; their raw responses stay in the working repository under
\texttt{tasks/20260928-comb0709-deeprecheck/}, not shipped in this deposit).
Summary by channel:

\begin{center}
\resizebox{\textwidth}{!}{%
\begin{tabular}{lll}
\hline
channel & queries & result \\
\hline
OEIS, sequence search \cite{oeis} & $10$ fresh strings (third round);
$50+$ independent strings for this sequence across rounds; $60+$ across the
three-candidate batch per the shared record, not shipped in this deposit &
zero hits for the main sequence: full string,
prepend-$a(0)$, prepend-zero, drop-first, drop-two, $\times 2$, mid-window
$6$-term and tail $7$-term windows, characteristic-polynomial coefficient
string, numerator string, mod-$3$ subsequences \\
OEIS, keyword & $5$ new strings (third round) & zero same-shape; four
groups of short-string generic hits, read entry by entry and excluded \\
OEIS, registry cross-read & $3$ records & A033506 signature
$(4,14,0,-10,0,1)$ read in full (the same-spectrum finding)
\cite{oeis-a033506}; A033507 \cite{oeis-a033507} and A030186
\cite{oeis-a030186} as mother-family anchors \\
zbMATH & $5$ & zero same-shape (two single hits read and excluded) \\
arXiv & $5$ & zero same-shape (two exact-phrase queries return total $0$;
one $49$-hit query all noise) \\
OpenAlex / Crossref & $10$ & zero same-shape across the alias set \\
Math Stack Exchange & $1$ & zero hits (engine-side query; native API
throttled during the review) \\
public Lean ecosystem & $4$ layers & zero: no matching-count theorem in
mathlib; compfiles zero; sequencelib structurally absent (no A-number);
GitHub code search on this sequence's large terms total $0$ \\
\hline
\end{tabular}%
}
\end{center}

The highest-risk neighbourhood was closed at the full-text level. Oh's
\emph{State matrix recursion method and monomer--dimer problem} \cite{oh2019}
contains, as its Theorem~4, the general formula for coverings with a fixed
monomer set --- the closest general method in the literature to this note's
object. The statistic is different in shape: Oh's $g_{m\times n}(S)$ counts
coverings whose monomers are \emph{exactly} the fixed set $S$, whereas $a(n)$
counts \emph{all} matchings of the graph obtained by \emph{deleting} the
side-middle cell (any number of monomers). The full text was read section by
section (ar5iv rendering): it contains no instance table, no OEIS reference,
and no recurrence for this sequence, and a literal search for twenty-three of
the large terms returned a single trivial substring hit. The three nearest
references in its bibliography were read in their arXiv versions --- Kong
\cite{kong}, Tzeng and Wu \cite{tzeng-wu}, Wu \cite{wu2006} --- and all three
concern the perfect-matching side with a single, unpositioned vacancy, a
different statistic again; the twelve citing papers of \cite{oh2019} found in
the review contain none on defect-grid full matching counts. The
neighbourhood is therefore recorded as a method-level generality, not an
occupying record for this sequence.

One correction produced by the third round is part of the record and is
load-bearing for this note's claims: the recurrence has the same
characteristic polynomial as the already-registered A033506
\cite{oeis-a033506}, signature $(4,14,0,-10,0,1)$ verbatim, so the earlier
rounds' ``new sequence, new recurrence'' wording was tightened to ``new
sequence'' --- the recurrence's order and coefficients are not new, only the
initial values, that is, the sequence itself. This note's wording follows the
corrected tier, and the correction is the reason no novelty claim is attached
to the recurrence anywhere in this note. The two original-language records
that fix this wording are quoted verbatim below.

\begin{lstlisting}[caption={Verbatim quote, original record (Chinese, kept untranslated): the priority fence locked on the source candidate card and repeated in the frozen statement file header (\texttt{proofs/01-comb08-sidemid-statements.lean}, lines 18--20).}, basicstyle=\ttfamily\footnotesize,frame=none,label={lst:fence}]
占位围栏：价值级仅「新序列」——六阶递推与 A033506（无缺陷 3×n 全匹配）同特征多项式（签名 (4,14,0,−10,0,1) 逐字相同），禁止对递推阶与系数作新颖性主张；PM≡0 为奇偶+洞色双重否决的真空命题，不构成占位亦不作价值主体。
\end{lstlisting}

\begin{lstlisting}[caption={Verbatim quote, original record (Chinese, kept untranslated): the C9 third-round verdict, from \texttt{audit/c9-record.md} (source \texttt{deep-recheck-combm-02.md}, section 0 and section 6).}, basicstyle=\ttfamily\footnotesize,frame=none,label={lst:c9verdict}]
**C9 四态：查无占位（三轮一致，本审支持前两审结论）。**
**价值分级：新序列——「新递推」主张经本审机器判定不成立（递推与已挂名的 A033506 同特征多项式），这是对前两审分级措辞的修正（收紧，不影响四态裁决）。**
- **C9 四态：查无占位**（证据链同共享记录 §8）。
- **价值分级：新序列**（递推非新——与 A033506 同谱；「新序列·新递推」措辞按本审修正）。派生 PM≡0 为经典着色阻塞实例，不构成占位亦不作价值主体。
\end{lstlisting}

The perfect-matching side of the neighbourhood is already occupied for
neighbouring defect configurations: the even-$n$ perfect-matching subsequence
of one neighbouring defect is A061278 \cite{oeis-a061278}, and the
perfect-matching subsequence of another neighbouring defect configuration is
A129113 \cite{oeis-a129113}. For this note's configuration the
perfect-matching count is identically zero, so there is no subsequence to
occupy and nothing to claim; the two registry entries are reported as
context, not as prior art for any statement of this note.

Given that evidence this note claims what the review supports: the sequence
was not found to be occupied. The novelty of Theorems~\ref{thm:interleave},
\ref{thm:mod2} and \ref{thm:mod4} is consequently negative-search evidence,
not a first-discovery assertion, and should be read as such. Disclosed holes:
the general web-search layer was unreachable in the review environment, so
the literature evidence rests on OpenAlex, zbMATH, arXiv, Crossref and
engine-side queries; the mother-family monographs \cite{lundow1996} were
identified through the registry records but their full texts were not read;
one paywalled reference was not read; the Math Stack Exchange native API was
throttled; and sequencelib is structurally absent for sequences without an
A-number. If those channels later return a hit on this sequence, the
positioning in this section weakens accordingly; the theorems themselves do
not depend on it.

\section{Scope and limitations}

Grade and boundary, stated as the audit states them. (a) The theorems cover
the recurrence-defined sequences only; their equality with the tiling counts
of the side-defected grid is not kernel-proved and is
Conjecture~\ref{conj:bridge}, whose empirical support is enumeration by exact
programs (three structurally different methods; seeds $n\le 14$ reproduced
digit-for-digit; deep window $n\le 50$; the enumerators themselves are not
verified, and the scripts and their outputs stay in the working repository
under \texttt{tasks/20260928-comb08-sidemid/}). (b) The recurrence's order and
coefficients are not new: its characteristic polynomial is shared with the
registered A033506 \cite{oeis-a033506}, signature $(4,14,0,-10,0,1)$ verbatim,
and no novelty is claimed for them anywhere in this note; what is claimed is
the sequence itself, at the value tier \emph{new sequence}. (c) The
perfect-matching count of this configuration is identically zero --- a vacuum
proposition forced by parity and by the classical colouring obstruction ---
and is reported as neither an occupying record nor a value subject. (d) The
gate-four sign-off of the underlying task was exercised by the author's
instruction to deposit this claim on 2026-09-28, while the task report file
still carries its draft header reading, in the original, ``awaiting manual
sign-off at gate four''; both facts are recorded here as they stand, and the
\texttt{tasks/} products were not modified. (e) This note carries no
literature review and is a deposit record, not a survey; the narrative paper
for the same result is separate.

\section*{Data availability}

The deposit package for this claim contains the claim note (LaTeX source and
PDF), the Lean sources (\texttt{proofs/}: frozen statement file, final proof
file, \texttt{lean-toolchain}), the verification and audit logs, and the full
C9 query log with zero-hit entries (\texttt{audit/}). Zenodo: the DOI is
reserved at upload by the author; that field is currently unfilled and is
written \texttt{RESERVED-DOI} in this note, to be replaced by the reserved
DOI in this note, in \texttt{metadata/README.md} of this package and in a
recompiled PDF at upload time; the manual steps are listed in
\texttt{UPLOAD.md} in the source repository (not shipped in this deposit).
Public mirror of this deposit: the claim lane's repository
\url{https://github.com/Aurora0134/ai4math-claims}, directory
\texttt{grid3n-sidemid-notch/} (pushed under the author's instruction of
2026-09-28; the commit and timestamp are recorded in the source repository's
\texttt{claims/grid3n-sidemid-notch/card.md}). Claim note under CC BY 4.0,
code under MIT.

\section*{Use of generative AI}

(a) AI performed: candidate generation and the exclusivity review,
autoformalisation of the three statements, the roundtrip semantic check,
proof search and proof-script writing, error classification and repair, and
the assembly of this note including its verbatim listings. (b) Model, node
and budget: all task-stage sampling ran on one node, wire-id
\texttt{anthropic/a6api-main/kimi-k3} (the task card's snapshot; a same-day
drift observation to \texttt{anthropic/stepfun/step-5-preview} is disclosed
in the source task's ledger rather than replaced), with no cross-node
switching, and consumed $14$ agent-run samplings of a permitted $20$, $1$
\texttt{llm-call} of a permitted $6$, $26$ kernel compiles of a permitted
$32$ and $9$ LaTeX rounds of a permitted $10$, copied verbatim from the task
ledger's final reconciliation. This note's assembly consumed $0$ pipeline
LLM calls against a cap of $4$ and was dispatched as one agent-run sampling
on the current session node (\texttt{stepfun/step-5-preview}). (c) Final
adjudication: all statements and proofs reported here were checked by the
Lean 4 kernel: every proof compiles with no \texttt{sorry}, and
\texttt{\#print axioms} reports only \{propext, Quot.sound\}, a subset of the
standard whitelist \{propext, Quot.sound, Classical.choice\}. (d) The AI
system is not an author; the human author reviewed the evidence and is
responsible for all content.

\section*{Author contributions}

CRediT-style: the human author adjudicated topic selection, statement
semantics and the deposit decision; the AI4Math pipeline (an LLM) generated
the candidate propositions, the proof searches and the text drafts. The AI
system is not listed as an author.

\begin{thebibliography}{99}
% Reused from papers/grid3n-sidemid-notch/paper.tex (source repository, not
% shipped in this deposit), where every entry below was verified on 2026-09-28
% before being typed in: Crossref REST metadata for the three DOIs (title
% matched exactly), arXiv export API title matches for the four preprints, OEIS
% fmt=json full-record reads plus landing-page HTTP 200 for the five sequence
% entries, and a GitHub API tag read for the mathlib revision. Raw responses:
% papers/grid3n-sidemid-notch/audit/bib-verify/ (source repository, not shipped
% in this deposit). No entry is quoted from memory; an HTTP 200 is not
% existence evidence, title metadata has to be matched (the template's _27 Lean
% DOI, which resolves to a different paper, is not used here).
\bibitem{lean4}
L.~de Moura and S.~Ullrich.
\newblock The Lean 4 theorem prover and programming language.
\newblock \emph{Automated Deduction -- CADE-28}, LNCS, pp.~625--635, 2021.
\url{https://doi.org/10.1007/978-3-030-79876-5_37}

\bibitem{mathlib}
The mathlib Community.
\newblock The Lean mathematical library.
\newblock \emph{CPP 2020: Proceedings of the 9th ACM SIGPLAN International
Conference on Certified Programs and Proofs}, pp.~367--381, 2020.
\url{https://doi.org/10.1145/3372885.3373824}

\bibitem{mathlib-v434}
The mathlib Community.
\newblock mathlib4, tag \texttt{v4.34.0} (revision
\texttt{5ed29652}\allowbreak\texttt{56430c36}\allowbreak\texttt{49e86755}\allowbreak\texttt{f9576b54}\allowbreak\texttt{eca72435}).
\newblock \url{https://github.com/leanprover-community/mathlib4}

\bibitem{oh2019}
S.~Oh.
\newblock State matrix recursion method and monomer--dimer problem.
\newblock \emph{arXiv:1901.07847 [math.CO]}, 2019.
\url{https://arxiv.org/abs/1901.07847}

\bibitem{kong}
Y.~Kong.
\newblock Packing dimers on $(2p+1)\times(2q+1)$ lattices.
\newblock \emph{arXiv:1410.8059 [cond-mat.stat-mech]}, 2014.
\url{https://arxiv.org/abs/1410.8059}

\bibitem{tzeng-wu}
W.-J.~Tzeng and F.~Y.~Wu.
\newblock Dimers on a simple-quartic net with a vacancy.
\newblock \emph{arXiv:cond-mat/0203149}, 2002.
\url{https://arxiv.org/abs/cond-mat/0203149}

\bibitem{wu2006}
F.~Y.~Wu.
\newblock The Pfaffian solution of a dimer-monomer problem: single monomer on
the boundary.
\newblock \emph{arXiv:cond-mat/0607647}, 2006.
\url{https://arxiv.org/abs/cond-mat/0607647}

\bibitem{godsil-gutman}
C.~D. Godsil and I.~Gutman.
\newblock On the theory of the matching polynomial.
\newblock \emph{Journal of Graph Theory} 5 (1981), pp.~137--144.
\url{https://doi.org/10.1002/jgt.3190050203}

\bibitem{lundow1996}
P.~H. Lundow.
\newblock Computation of matching polynomials and the number of 1-factors in
polygraphs.
\newblock Research reports No.~12, Department of Mathematics, Ume{\aa}
University, 1996; as referenced in OEIS A033506 \cite{oeis-a033506}.

\bibitem{oeis}
The OEIS Foundation.
\newblock \emph{The On-Line Encyclopedia of Integer Sequences}.
\newblock \url{https://oeis.org/} --- searched for the sequences and their
variant strings, with no entry found for the main sequence (see
Section~\ref{sec:novelty}).

\bibitem{oeis-a033506}
The OEIS Foundation.
\newblock Entry A033506 (P.~H. Lundow): number of matchings in graph
$P_3\times P_n$; index entry for linear recurrences with constant
coefficients, signature $(4,14,0,-10,0,1)$.
\newblock \emph{The On-Line Encyclopedia of Integer Sequences}.
\url{https://oeis.org/A033506}

\bibitem{oeis-a033507}
The OEIS Foundation.
\newblock Entry A033507: number of matchings in graph $P_4\times P_n$; index
entry signature $(9,41,-41,-111,91,29,-23,-1,1)$.
\newblock \emph{The On-Line Encyclopedia of Integer Sequences}.
\url{https://oeis.org/A033507}

\bibitem{oeis-a030186}
The OEIS Foundation.
\newblock Entry A030186: $a(n)=3a(n-1)+a(n-2)-a(n-3)$, the $2\times n$ grid
mother-family anchor.
\newblock \emph{The On-Line Encyclopedia of Integer Sequences}.
\url{https://oeis.org/A030186}

\bibitem{oeis-a061278}
The OEIS Foundation.
\newblock Entry A061278: $a(n)=5a(n-1)-5a(n-2)+a(n-3)$; the perfect-matching
subsequence of a neighbouring defect configuration (context only, not prior
art for any statement of this note).
\newblock \emph{The On-Line Encyclopedia of Integer Sequences}.
\url{https://oeis.org/A061278}

\bibitem{oeis-a129113}
The OEIS Foundation.
\newblock Entry A129113: expansion of
$x^3/(1-x-5x^2-x^3+x^4)$; the perfect-matching subsequence of another
neighbouring defect configuration (context only, not prior art for any
statement of this note).
\newblock \emph{The On-Line Encyclopedia of Integer Sequences}.
\url{https://oeis.org/A129113}
\end{thebibliography}

\end{document}
